\section{Related work}
\label{sec:related}

\paragraph{Cooperative web conventions and honor-based machine directives.}
The closest conceptual ancestor of Recuse is the Robots Exclusion Protocol, the
\texttt{robots.txt} convention introduced by Martijn Koster in 1994 and only
recently codified as an Internet standard in RFC~9309
\citep{rfc9309}. The defining property of this lineage is that it is
\emph{voluntary and honor-based}: the server publishes a machine-readable
directive describing which automated clients may access which resources, and
compliant crawlers are expected to withdraw of their own accord. RFC~9309
standardized parsing, error handling, and caching semantics but deliberately did
not introduce any enforcement mechanism---compliance remains a matter of the
client's good faith. The same honor-based philosophy animates ongoing
standardization of \emph{AI-usage} preferences: the IETF AI Preferences (aipref)
working group is extending this model to signal whether content may be used for
AI training or inference \citep{aipref-attach}. Recuse follows directly in this
tradition but differs in target and timing: rather than a crawl-time directive
about bulk content harvesting, it defines an \emph{in-band, per-request
deny-signal} emitted by a live server so that an LLM \emph{agent}---not a
crawler---voluntarily recuses itself from a specific resource at access time.

\paragraph{LLM agent access control via gateways, tools, and authorization.}
A growing body of practice mediates agent access through external authorization
layers. The Model Context Protocol (MCP) standardizes how LLM applications
connect to tools and data sources, and its security model layers OAuth~2.1-based
authorization and explicit user-consent flows on top of tool invocation
\citep{mcp2025}. Notably, the MCP specification itself acknowledges that it
``cannot enforce these security principles at the protocol level,'' delegating
consent and access control to host implementations. This captures a structural
feature of the gateway/proxy approach generally: authorization lives
\emph{external to the resource}, in a mediating client, gateway, or token issuer
that decides whether a call is permitted before it reaches the server. Recuse is
complementary but architecturally inverted---the signal originates \emph{at the
resource itself}, is legible to the agent, and asks the agent to decline, rather
than asking an intermediary to block.

\paragraph{Relationship- and role-based authorization.}
Modern fine-grained authorization is dominated by relationship-based access
control (ReBAC), popularized by Google's Zanzibar, which stores
object--relation--user tuples and evaluates authorization as graph reachability
at massive scale \citep{pang2019zanzibar}. Its open-source descendant OpenFGA
brings the same model to general developers \citep{openfga}. These systems answer
the question ``is principal $X$ permitted to act on object $Y$?'' and they do so
as an \emph{external} policy decision point consulted by the application. Recuse
addresses a different question entirely: it is not a permission-evaluation engine
but a \emph{published deny-signal} the resource volunteers to a cooperating
agent. Crucially, our measurements show that an on-host policy signal can outrank
prompt-level authorization in agent behavior---a finding that ReBAC/RBAC systems,
which sit outside the resource and outside the model's instruction context, are
not designed to address.

\paragraph{Instruction hierarchy and authority conflicts in LLMs.}
When a server's deny-signal contradicts a user or system prompt that authorizes
access, the agent faces an authority conflict. \citet{wallace2024instruction}
formalize this as an \emph{instruction hierarchy}, proposing training that teaches
models to prioritize privileged instructions and selectively ignore
lower-privilege ones. Our work is empirically adjacent: we find that a Recuse
signal delivered in-band can, in practice, outrank explicit prompt authorization,
which both validates the premise that models arbitrate among competing
authorities and surfaces a governance lever---the resource's own voice---that the
instruction-hierarchy literature does not consider as a distinct source of
authority.

\paragraph{Prompt injection and the confused-deputy problem.}
A large security literature treats untrusted content reaching an LLM agent as an
attack surface. \citet{greshake2023injection} introduced \emph{indirect prompt
injection}, showing that adversarial instructions embedded in retrieved data can
hijack real-world LLM-integrated applications; structurally this is a modern
instance of the classic \emph{confused deputy} \citep{hardy1988confused}, in
which a privileged agent is tricked into misusing its authority on behalf of a
less-privileged party. This framing is important for honestly positioning Recuse:
because a deny-signal is just more in-band text, a malicious server could in
principle emit it to manipulate an agent, and a malicious client can trivially
ignore it. We therefore explicitly frame Recuse as a \emph{cooperative governance
signal, not a security control}---it presumes a compliant agent and defers
adversarial enforcement, exactly the threat model the prompt-injection literature
targets.

\paragraph{Machine-readable consent for agents, and the measurement gap.}
The idea of a \texttt{robots.txt}-style, machine-readable policy aimed
specifically at autonomous web agents is only just emerging. Most directly, in recent prior
work \citet{marro2026permission} propose \emph{permission manifests} (an
\texttt{agent-permissions.json} file) that let sites declare which agent
interactions are permitted, echoing our motivation; that work proposes the
manifest mechanism and format but does not report an empirical measurement of
whether deployed agents actually comply. Recuse differs along two axes that
constitute our core contribution. First, mechanism: we standardize an
\emph{in-band, per-request, agent-legible deny-signal} emitted by the live server
during access, rather than a separately fetched static manifest. Second, and more
importantly, to the best of our knowledge \emph{we are the first to empirically
measure whether deployed LLM agents actually honor such a signal}. We are aware of
no prior work---neither the cooperative-convention standards
\citep{rfc9309,aipref-attach}, the agent-permission proposals
\citep{marro2026permission}, nor the authorization systems above---that reports a
controlled measurement of agent \emph{compliance} with a voluntary in-band
withdrawal request; this measurement gap is a central part of what we contribute.
We deliberately set aside the \emph{enforcement layer}---behavioral bot-management
and automation detection that distinguish automated from human traffic via
traffic patterns and biometric/fingerprint signals \citep{iliou2021bots}---as
orthogonal and complementary future work for the adversarial regime.
